\capitulo{7}{Líneas de trabajo futuras}
 Aunque la aplicación UV-Dermis ya ofrece  funcionalidades relevantes para la prevención del \textit{melanoma} cutáneo, existen múltiples posibilidades de mejora técnica y expansión futura. Como líneas de trabajo futuras para este proyecto se proponen las siguientes:
 \begin{itemize}
     \item \textbf{Geolocalización automática del usuario:} actualmente la app dispone de una selección manual de la zona geográfica del usuario, la incorporación de esta mejora permitiría un acceso y recomendación más rápido. Esta funcionalidad está ampliamente implementada en aplicaciones web y mejora significativamente la experiencia de usuario. Desde un punto de vista técnico integraríamos un nuevo lenguaje de programación, \textit{JavaScript}, mediante un paquete de \textit{R} llamado \textit{shinyjs} \cite{Shinyjs} que nos permitiría implementar nuevas acciones en nuestra app sin necesidad de saber en profundidad de este lenguaje. Se incluye este lenguaje nuevo ya que \textit{R} no tiene acceso a la API de geolocalización del navegador, de esta manera se produce un cambio de información una ver devuelta la respuesta de las coordenadas del usuario y se traducirían directamente a la selección automática de la provincia.
     \item \textbf{Incorporación de variables atmosféricas adicionales}: el modelo actual de la aplicación UV-Dermis se basa en variables ambientales clave como la radiación ultravioleta (UV), la temperatura máxima y la altitud por provincia. Sin embargo, diversos estudios científicos y organismos internacionales \cite{organizationGlobalSolarUV2002} reconocen que existen otros factores atmosféricos modificadores de la exposición real a la radiación UV que incide sobre la piel. Entre los más relevantes destacan la nubosidad y la capa de ozono. La nubosidad modula cuánta radiación UV realmente llega al suelo, algunas nubes reducen la radiación, pero otras la intensifican \cite{organizationGlobalSolarUV2002}. Por otro lado, la capa de ozono actúa como escudo natural \cite{UltravioletRadiationHow2001}: cuando está más delgada, se filtra menos radiación UVB, aumentando el riesgo para la piel. Su incorporación permitiría a la aplicación ofrecer recomendaciones más precisas dado a que se introducen factores determinantes en la lógica. Además, los datos de ambas variables están disponibles en la AEMET.
     \item \textbf{Implementación de un sistema de notificaciones automáticas:} actualmente, UV-Dermis ofrece recomendaciones tras una consulta activa por parte del usuario. Una solucion a esto es vía email, técnica ampliamente extendida y accesible para la gran mayoría de los usuarios ya que hoy en día todos poseemos un correo electrónico. Esta nueva incorporación se implementaría con la ayuda del paquete de \textit{R} \textit{emayili} que permite construir y enviar correos electrónicos desde \textit{R} empleando servidores SMTP\footnote{SMTP significa \textit{Simple Mail Transfer Protocol} o Protocolo Simple de Transferencia de Correo. Es un protocolo de comunicación que se usa para enviar y recibir correos electrónicos a través de internet.}\cite{collierDatawookieEmayili2025}. Con esta mejora transformaría UV-Dermis en una app proactiva y no solo consultiva, aumentando el impacto preventivo en la población interesada.
     Por otra, parte esta mejora deberá cumplir con el Reglamento General de Protección de Datos (RGPD). Esto implica obtener el consentimiento explícito e informado del usuario, asegurar la encriptación de datos personales y ofrecer opciones de baja voluntaria del sistema de alertas si el usuario lo solicita.
 \end{itemize}

